\documentclass[11pt]{article}
\usepackage[margin=1in]{geometry}
\usepackage{amsmath,amssymb,mathtools}
\usepackage{fontspec,unicode-math}
\setmainfont{Latin Modern Roman}
\setsansfont{Latin Modern Sans}
\setmathfont{Latin Modern Math}
\usepackage{xurl,hyperref}
\hypersetup{colorlinks=true,urlcolor=blue,linkcolor=blue}
\newcommand{\sref}[2]{\hyperlink{source:#1:#2}{[S#2]}}
\newcommand{\cataloguescope}[1]{\paragraph{Recorded mathematical question.}#1}
\begin{document}
\section{Multiplicative connection bounds above criticality}\label{statement}

\subsection{Definitions and mathematical statement}

\paragraph{Shared notation.}$\mathbb N=\{1,2,\ldots\}$, and $\mathbb Z,\mathbb Q,\mathbb R,\mathbb C$ denote the integers, rationals, reals and complex numbers. Any other conventions are specified locally.


Let $G$ be a connected locally finite nonamenable hyperbolic graph. Hyperbolic means its graph-metric geodesic triangles are uniformly thin. Nonamenable means its Cheeger constant $\inf_{0<|A|<\infty}|\partial_E A|/|A|$ is positive, where $\partial_E A$ is the set of edges with one endpoint in $A$. In Bernoulli bond percolation, each edge is independently open with probability $p$. Write $x\leftrightarrow y$ for an open connection and $p_c$ for the threshold for infinite clusters. A geodesic is a shortest unit-edge path.
\paragraph{Statement recorded in the supplied source.}
For which $p\in[0,1]$ is there $C=C(G,p)<\infty$ such that
\[\mathbb P_p(x\leftrightarrow y)\le C\mathbb P_p(x\leftrightarrow z)\mathbb P_p(z\leftrightarrow y)\]
for all vertices $x,y$ and every $z$ on a geodesic from $x$ to $y$? In particular, decide this at $p_c$ and for some $p>p_c$. \sref{statement}{2}
\subsection{Short English statement}
When does percolation connectivity on a nonamenable hyperbolic graph satisfy a uniform multiplicative upper bound through points on geodesics?
\subsection{Sources}
\begin{description}
\item[\hypertarget{source:statement:1}{S1}] ULAM.AI and the UnsolvedMath Contributors, \emph{UnsolvedMath: A Curated Collection of Open Mathematics Problems}, catalogue entry. CC BY 4.0. \url{https://huggingface.co/datasets/ulamai/UnsolvedMath}.
\item[\hypertarget{source:statement:2}{S2}] Itai Benjamini, \emph{Coarse Geometry and Randomness}, Saint-Flour lecture notes, version dated 30 October 2013; original author manuscript preserved in Arquivo.pt capture of 31 December 2020, Section 9.2, Open Problem 9.14 and immediately preceding paragraph, p. 66. \url{https://arquivo.pt/noFrame/replay/20201231041548id_/http://www.wisdom.weizmann.ac.il/~itai/stflouraug24.pdf}.
\end{description}
\label{end:statement}
\end{document}
